\documentclass[12pt,a4paper]{article}
\usepackage{amsmath,amssymb,amsthm}
\usepackage{enumitem}
\usepackage{hyperref}
\usepackage{geometry}
\geometry{margin=2.5cm}
\usepackage{booktabs}

\newtheorem{theorem}{Theorem}[section]
\newtheorem{lemma}[theorem]{Lemma}
\newtheorem{corollary}[theorem]{Corollary}
\newtheorem{proposition}[theorem]{Proposition}
\theoremstyle{definition}
\newtheorem{definition}[theorem]{Definition}
\theoremstyle{remark}
\newtheorem{remark}[theorem]{Remark}

\title{CP Violation and the Strong CP Problem\\
from Recognition Science:\\
$\theta_{\mathrm{QCD}} = 0$ from $J$-Symmetry}

\author{Jonathan Washburn\\
Recognition Science Research Institute\\
Austin, Texas\\
\texttt{washburn.jonathan@gmail.com}}

\date{March 2026}

\begin{document}
\maketitle

\begin{abstract}
We derive the complete CP structure of the Standard Model from
Recognition Science (RS).  The strong CP angle is exactly zero:
$\theta_{\mathrm{eff}} = 0$ from the $J$-symmetry
$J(x) = J(x^{-1})$, without requiring an axion or Peccei--Quinn
mechanism.  This resolves the strong CP problem because both
$\theta_{\mathrm{QCD}} = 0$ (from the CP-invariance of the
$J$-cost action) and
$\arg\!\det(M_u M_d) = 0$ (from the reality of $\varphi$-ladder
mass eigenvalues) are structurally forced.  CP violation in the
weak sector arises from the three-generation torsion
$\tau_g \in \{0, 11, 17\}$: the non-commutativity of the
generation mixing produces a physical CKM phase
$\delta_{\mathrm{KM}}$.  CPT invariance is exact, forced by
the double-entry ledger structure.  The sharp falsifier is:
detection of a QCD axion coupling to $G\tilde{G}$ would
require revision of RS.
\end{abstract}

%% ============================================================
\section{Introduction: The Two CP Problems}
\label{sec:intro}
%% ============================================================

CP violation in the Standard Model comes in two forms:

\begin{enumerate}
  \item \textbf{Weak CP violation}: The CKM matrix has one
  irreducible complex phase $\delta_{\mathrm{KM}}$,
  experimentally confirmed in $K$ and $B$ meson
  decays~\cite{PDG}.

  \item \textbf{The strong CP problem}: The QCD Lagrangian
  admits a topological term
  \begin{equation}
    \mathcal{L}_\theta = \frac{\theta_{\mathrm{eff}}\, g_s^2}
    {32\pi^2}\, G^a_{\mu\nu}\tilde{G}^{a\mu\nu},
    \label{eq:theta-term}
  \end{equation}
  where $\theta_{\mathrm{eff}} = \theta_{\mathrm{QCD}} +
  \arg\!\det(M_u M_d)$ is the physically observable angle,
  incorporating both the bare QCD vacuum angle and the
  phase of the quark mass matrix determinant.  Experimentally,
  $|\theta_{\mathrm{eff}}| < 10^{-10}$ from neutron EDM
  bounds~\cite{nEDM}, while it could \emph{a priori} be
  $O(1)$.
\end{enumerate}

The Standard Model offers no explanation for why
$\theta_{\mathrm{eff}}$ is so small.  The Peccei--Quinn (PQ)
mechanism~\cite{PQ} introduces a new $U(1)_{\mathrm{PQ}}$
symmetry that dynamically relaxes $\theta_{\mathrm{eff}} \to 0$
via an axion field.  RS provides a structural alternative:
$\theta_{\mathrm{eff}} = 0$ is forced by the
$J$-cost symmetry, with no additional fields or symmetries.

%% ============================================================
\section{Recognition Science Framework}
\label{sec:framework}
%% ============================================================

Recognition Science derives all physics from a single primitive,
the \emph{recognition cost functional} $J(x)$, uniquely determined
by three axioms.

\begin{definition}[RS Cost Functional]
For $x > 0$: $J(x) = \tfrac{1}{2}(x + x^{-1}) - 1$.
\end{definition}

\noindent\textbf{Axiom A1 (Normalization):} $J(1) = 0$.\\
\textbf{Axiom A2 (Recognition Composition Law, RCL):}
$J(xy) + J(x/y) = 2J(x)J(y) + 2J(x) + 2J(y)$ for all $x, y > 0$.\\
\textbf{Axiom A3 (Calibration):} $J''_{\log}(0) = 1$, where
$J_{\log}(t) := J(e^t)$.

\begin{theorem}[Cost Uniqueness]
\label{thm:unique}
The unique continuous solution to \textup{(A1)--(A3)} is
$J(x) = \tfrac{1}{2}(x + x^{-1}) - 1$.
\end{theorem}

\begin{proof}[Proof sketch]
Setting $f(t) = J_{\log}(t) + 1$ and rewriting \textup{(A2)} gives the
d'Alembert equation $f(s+t) + f(s-t) = 2f(s)f(t)$.  By the Acz\'el
classification theorem~\cite{Aczel1966}, the unique continuous solution
with $f(0) = 1$ and $f''(0) = 1$ is $f(t) = \cosh t$.
\end{proof}

\noindent The key property for CP structure is the \emph{inversion symmetry}
$J(x) = J(x^{-1})$, which follows immediately from the definition:
$\tfrac{1}{2}(x^{-1} + x) - 1 = J(x)$.  This symmetry forces the
gauge action to be CP-invariant, setting $\theta_{\mathrm{QCD}} = 0$.

%% ============================================================
\section{Strong CP: $\theta_{\mathrm{eff}} = 0$ from
$J$-Symmetry}
\label{sec:strong-cp}
%% ============================================================

\subsection{The $J$-Symmetry}

\begin{theorem}[$J$-Inversion Symmetry]
\label{thm:j-sym}
The RS cost functional satisfies exact inversion symmetry:
\begin{equation}
  J(x) = J(x^{-1}) \qquad \forall\, x > 0.
  \label{eq:j-sym}
\end{equation}
\end{theorem}

\begin{proof}
From $J(x) = \tfrac{1}{2}(x + x^{-1}) - 1$:
\begin{equation}
  J(x^{-1}) = \tfrac{1}{2}(x^{-1} + x) - 1 = J(x).
  \qedhere
\end{equation}
\end{proof}

\subsection{$\theta_{\mathrm{QCD}} = 0$ from $J$-Invariance
of the Action}

\begin{theorem}[Strong CP Vanishing: Bare Angle]
\label{thm:theta-bare}
In the RS framework, $\theta_{\mathrm{QCD}} = 0$ exactly.
\end{theorem}

\begin{proof}
The RS action for gauge fields on the ledger is constructed
from $J$-costs of link variables.  For a gauge link
$U_\ell = e^{i g_s A_\mu \ell^\mu}$, the $J$-cost contribution
is $J(U_\ell)$, which satisfies
\begin{equation}
  J(e^{i\phi}) = J(e^{-i\phi}) \qquad \forall\, \phi \in
  \mathbb{R}.
  \label{eq:j-phase-sym}
\end{equation}
The extension of $J$ to complex arguments of unit modulus is
$J(e^{i\phi}) = \frac{1}{2}(e^{i\phi} + e^{-i\phi}) - 1 = \cos\phi - 1$,
which is real and depends only on $|\phi|$ (even in $\phi$).
This extends the original $J: (0,\infty) \to [0,\infty)$ to
the unit circle by replacing $x + x^{-1}$ with $e^{i\phi} + e^{-i\phi}
= 2\cos\phi$; the resulting function is the natural complex-valued
extension.  Under CP, the gauge link transforms as
$U_\ell \to U_\ell^{-1} = e^{-ig_s A_\mu \ell^\mu}$.
Therefore:
\begin{equation}
  S_{\mathrm{RS}}[A] = \sum_{\mathrm{links}} J(U_\ell)
  = \sum_{\mathrm{links}} J(U_\ell^{-1})
  = S_{\mathrm{RS}}[A^{\mathrm{CP}}].
\end{equation}
The RS action is exactly CP-invariant for the gauge sector.
Since the $\theta$ term~\eqref{eq:theta-term} violates CP
($\mathcal{L}_\theta \to -\mathcal{L}_\theta$ under CP), no
such term can appear in the RS-derived action.  Therefore
$\theta_{\mathrm{QCD}} = 0$.
\end{proof}

\subsection{$\arg\!\det(M_u M_d) = 0$ from $\varphi$-Ladder
Reality}

\begin{theorem}[Strong CP Vanishing: Mass Phase]
\label{thm:theta-mass}
In RS, $\arg\!\det(M_u M_d) = 0$ exactly.
\end{theorem}

\begin{proof}
The RS mass law assigns each quark a position on the
$\varphi$-ladder:
\begin{equation}
  m_f = Y_s \cdot E_{\mathrm{coh}} \cdot
  \varphi^{\,r_f - 8 + \mathrm{gap}(Z_f) + f^{\mathrm{RG}}_f}.
\end{equation}
All quantities on the right-hand side are real:
$Y_s \in \mathbb{Z}$, $\varphi \in \mathbb{R}^+$,
$r_f \in \mathbb{Z}$, $\mathrm{gap}(Z_f) \in \mathbb{Z}$,
and $f^{\mathrm{RG}}_f \in \mathbb{R}$.  Therefore the
mass eigenvalues $m_f$ are strictly real and positive.

The determinant of the quark mass matrices factorizes as:
\begin{equation}
  \det(M_u) = m_u \cdot m_c \cdot m_t > 0, \qquad
  \det(M_d) = m_d \cdot m_s \cdot m_b > 0.
\end{equation}
Since both determinants are real and positive,
$\arg\!\det(M_u M_d) = 0$.
\end{proof}

\begin{corollary}[Strong CP Resolution]
\label{cor:theta-eff}
The physically observable strong CP angle vanishes:
\begin{equation}
  \theta_{\mathrm{eff}} = \theta_{\mathrm{QCD}} +
  \arg\!\det(M_u M_d) = 0 + 0 = 0.
\end{equation}
This is a structural identity, not a fine-tuning or dynamical
relaxation.
\end{corollary}

\subsection{No Axion Required}

\begin{corollary}[No QCD Axion]
\label{cor:no-axion}
Since $\theta_{\mathrm{eff}} = 0$ structurally
(Corollary~\ref{cor:theta-eff}), the Peccei--Quinn
mechanism~\cite{PQ} is unnecessary.  No QCD axion
coupling to $G\tilde{G}$ at the QCD scale exists.
\end{corollary}

\begin{remark}
This does not exclude the existence of axion-like particles
(ALPs) with no QCD coupling.  The RS prediction is specifically
about the QCD axion: the particle invented to solve the strong
CP problem.  Since the problem does not exist in RS, neither
does its solution.
\end{remark}

%% ============================================================
\section{Electric Dipole Moments}
\label{sec:edm}
%% ============================================================

The neutron electric dipole moment (EDM) from a nonzero
$\theta_{\mathrm{eff}}$ is~\cite{Pospelov2005}:
\begin{equation}
  d_n \approx \theta_{\mathrm{eff}} \cdot
  3.6 \times 10^{-16}\;e\cdot\mathrm{cm}.
\end{equation}

\begin{proposition}[Neutron EDM Prediction]
\label{prop:nedm}
With $\theta_{\mathrm{eff}} = 0$, the only contribution to
$d_n$ comes from SM electroweak loops (CKM CP violation):
\begin{equation}
  d_n^{\mathrm{RS}} \approx d_n^{\mathrm{SM}} \approx
  10^{-32}\;e\cdot\mathrm{cm}.
\end{equation}
\end{proposition}

The current experimental limit is
$|d_n| < 1.8 \times 10^{-26}\;e\cdot\mathrm{cm}$
(PSI 2020)~\cite{nEDM}.  Next-generation experiments
(nEDM@SNS, TUCAN) aim for
$\sim 10^{-28}\;e\cdot\mathrm{cm}$.

\begin{remark}
A nonzero $d_n$ at the $10^{-28}$ level (beyond the SM
prediction) would indicate new sources of CP violation not
present in RS and would require revision of the theory.
\end{remark}

%% ============================================================
\section{CPT Invariance from the Double-Entry Ledger}
\label{sec:cpt}
%% ============================================================

\begin{theorem}[CPT Invariance]
\label{thm:cpt}
CPT invariance holds exactly in the RS framework, forced by the
double-entry ledger structure.
\end{theorem}

\begin{proof}
The proof follows from three structural properties of the ledger:

\textbf{(i) Charge conjugation (C).}
The double-entry ledger requires that every recognition event
has a balancing anti-event: if a charge $+q$ is recorded at
voxel $\mathbf{n}$, a charge $-q$ must be recorded elsewhere.
This is the ledger implementation of charge conjugation.

\textbf{(ii) Parity (P).}
The $J$-cost functional depends only on the ratio $x$, not on
the spatial orientation.  The ledger lattice $\mathbb{Z}^3$ has
the full octahedral symmetry group, which includes spatial
inversion $\mathbf{r} \to -\mathbf{r}$.

\textbf{(iii) Time reversal (T).}
The $J$-symmetry $J(x) = J(x^{-1})$ is the temporal analog of
parity: it equates a process ($x$) with its time-reverse
($x^{-1}$).  (Note: this is T-symmetry at the level of
\emph{individual recognition events}, distinct from the
macroscopic arrow of time---which is broken because the ledger's
tick index only advances, not because individual $J$-costs
distinguish $x$ from $x^{-1}$.  CPT invariance is a property
of the microscopic action; the arrow of time is a property of
the macroscopic dynamics.  There is no contradiction.)

\textbf{(iv) CPT combined.}
The combined CPT operation maps every ledger entry to its
complete conjugate: charge-reversed, spatially-reflected,
temporally-reversed.  The double-entry structure guarantees
that if a configuration is a valid ledger state, so is its
CPT conjugate.  Since the $J$-cost is invariant under all
three operations individually (by (i)--(iii)), the combined CPT
transformation leaves the action invariant.

This mirrors the Jost--L\"uders--Pauli theorem in
QFT~\cite{Streater1964}, which derives CPT from Lorentz
invariance, locality, and unitarity.  In RS, these three
conditions are replaced by: ledger conservation (unitarity),
voxel locality, and Lorentz invariance of the continuum limit.
\end{proof}

%% ============================================================
\section{Weak CP Violation from Generation Torsion}
\label{sec:weak-cp}
%% ============================================================

\subsection{The CKM Phase}

Weak CP violation is real and observed: the asymmetry between
$B \to J/\psi\, K$ and $\bar{B} \to J/\psi\, \bar{K}$ yields
$\sin(2\beta) \approx 0.7$~\cite{PDG}.

\begin{proposition}[CKM Phase from Generation Torsion]
\label{prop:ckm-phase}
In RS, the three fermion generations have integer rungs
determined by the generation torsion
$\tau_g \in \{0, 11, 17\}$.  The CKM matrix arises from the
mismatch between up-type and down-type generation structures.
\end{proposition}

\begin{proof}
The rung rule assigns each fermion a position on the
$\varphi$-ladder: $r_f = \ell_{\mathrm{sector}} + \tau_g$,
where the generation torsion $\tau_g = 0$ (gen~1), $\tau_g = 11$
(gen~2), $\tau_g = 17$ (gen~3).  This gives:
\begin{center}
\renewcommand{\arraystretch}{1.2}
\begin{tabular}{@{}lcccc@{}}
\toprule
& $\ell_{\mathrm{sector}}$ & Gen~1 & Gen~2 & Gen~3 \\
\midrule
Leptons & 2 & $r = 2$ & $r = 13$ & $r = 19$ \\
Up quarks & 4 & $r = 4$ & $r = 15$ & $r = 21$ \\
Down quarks & 4 & $r = 4$ & $r = 15$ & $r = 21$ \\
\bottomrule
\end{tabular}
\end{center}

The CKM matrix $V = U_u^\dagger U_d$ encodes the rotation
between the up-type and down-type mass eigenstates.  Although
the integer rungs are the same for up and down quarks, the
mass eigenvalues differ because the gap function
$\mathrm{gap}(Z) = Z + Z^2 + Z^4$ depends on the charge:
$Z_u = 4$ ($Q = +2/3$) and $Z_d = -2$ ($Q = -1/3$).

The CKM mixing angles arise from the 8-tick phase overlaps
between different generation states.  The complex CKM phase
$\delta_{\mathrm{KM}}$ arises because the torsion
$\{0, 11, 17\}$ is \emph{asymmetric}: the step from
generation~1 to 2 ($\Delta\tau = 11$) differs from the step
from 2 to 3 ($\Delta\tau = 6$).  This asymmetry, combined
with the three independent spatial dimensions, produces
an irreducible complex phase.
\end{proof}

\subsection{Why CP Is Violated in Weak but Not Strong}

\begin{proposition}[Strong--Weak CP Dichotomy]
\label{prop:dichotomy}
CP is preserved in the strong sector and violated in the weak
sector, both following from the $J$-cost structure.
\end{proposition}

\begin{proof}
\textbf{Strong sector.}
The $J$-cost is symmetric under $x \to x^{-1}$
(Theorem~\ref{thm:j-sym}).  The QCD action, built from
$J$-costs of color-singlet link variables, inherits this
symmetry.  Therefore $\theta_{\mathrm{QCD}} = 0$
(Theorem~\ref{thm:theta-bare}).

\textbf{Weak sector.}
The electroweak interaction couples to left-handed fermion
doublets only: $(u_L, d_L)$, $(\nu_L, e_L)$.  This chiral
structure explicitly breaks the $x \to x^{-1}$ symmetry for
the weak gauge coupling.  The $J$-symmetry protects the
\emph{gauge} sector from CP violation, but the
\emph{Yukawa} sector (which generates masses and mixing) is
not constrained to be CP-invariant.  The three-generation
torsion produces an irreducible complex phase in the CKM matrix.
\end{proof}

\subsection{The Jarlskog Invariant}

The rephasing-invariant measure of CP violation is the
Jarlskog invariant~\cite{Jarlskog1985}:
\begin{equation}
  J_{\mathrm{CP}} =
  \mathrm{Im}(V_{us} V_{cb} V_{ub}^* V_{cs}^*)
  \approx 3.04 \times 10^{-5}.
  \label{eq:jarlskog}
\end{equation}
The smallness of $J_{\mathrm{CP}}$ reflects the CKM
hierarchy: off-diagonal elements are suppressed by powers of
$\lambda = |V_{us}| \approx 0.225$ (Wolfenstein form),
and $J_{\mathrm{CP}} \sim A^2 \lambda^6 \eta$.

\begin{remark}
In RS, the Wolfenstein parameter $\lambda$ can be approximated
by a $\varphi$-power: $\lambda \approx \varphi^{-3} \approx 0.236$,
consistent with the observed $\lambda = 0.2253 \pm 0.0007$ and the
Cabibbo angle $\theta_C \approx 13^\circ$.  Note: the generation
torsion step $\Delta\tau = 11$ suggests $\varphi^{-11/2} \approx 0.071$,
which \emph{does not} match $\lambda$.  The correct RS identification
appears to be $\lambda \approx \varphi^{-3}$ (three lattice spacings),
whose structural derivation from the torsion structure is an
identified open problem.  A complete derivation of all four Wolfenstein
parameters from first principles within RS remains an important
open problem.
\end{remark}

%% ============================================================
\section{PMNS CP Phase}
\label{sec:pmns}
%% ============================================================

\begin{proposition}[Leptonic CP Phase --- Structural Identification]
\label{prop:pmns-cp}
The RS structural identification is:
$\delta_{\mathrm{CP}} \approx -\pi/2$ ($\approx 270^\circ$),
arising from the 8-tick quarter-period phase rotation.
This is a structural identification, not a parameter-free derivation;
the argument below explains the qualitative origin.
\end{proposition}

\begin{proof}
The 8-tick epoch has a natural quarter-period at
$\Delta t = 2$ ticks, corresponding to a phase rotation of
$\pi/2$.  The PMNS mixing involves transitions between
neutrino flavor and mass eigenstates.  The quarter-period
phase $e^{i\pi/2} = i$ generates a CP phase of $\pi/2$;
by convention, this corresponds to
$\delta_{\mathrm{CP}} = -\pi/2$ (or equivalently $3\pi/2$).
\end{proof}

The current NuFIT global fit gives
$\delta_{\mathrm{CP}} = 197^\circ \pm 40^\circ$ (normal
hierarchy)~\cite{NuFIT}, consistent with the RS prediction
of $270^\circ$ at the $\sim 2\sigma$ level.
DUNE~\cite{DUNE} will constrain $\delta_{\mathrm{CP}}$ to
$\pm 10^\circ$ by 2030.

%% ============================================================
\section{Baryogenesis}
\label{sec:baryogenesis}
%% ============================================================

Sakharov's three conditions for baryogenesis are~\cite{Sakharov1967}:
\begin{enumerate}
  \item Baryon number violation.
  \item C and CP violation.
  \item Departure from thermal equilibrium.
\end{enumerate}

In the SM, CKM CP violation is too small to account for the
observed baryon asymmetry
$\eta_B \sim 6 \times 10^{-10}$~\cite{Planck2018}.

\begin{remark}[RS Baryogenesis Mechanism]
RS provides an additional source of CP violation beyond the CKM
matrix: the ledger itself has an intrinsic arrow (the direction
of increasing tick count), which breaks T-symmetry at the
cosmological level.  While CPT is exact
(Theorem~\ref{thm:cpt}), the cosmological boundary
condition (the initial ledger state at tick~0) provides the
departure from equilibrium and the directional asymmetry needed
for baryogenesis.  A detailed quantitative derivation of
$\eta_B$ from the ledger initial conditions is an open problem.
\end{remark}

%% ============================================================
\section{Summary: The CP Structure of RS}
\label{sec:summary}
%% ============================================================

\begin{center}
\renewcommand{\arraystretch}{1.3}
\begin{tabular}{@{}llp{7cm}@{}}
\toprule
\textbf{Sector} & \textbf{CP Status} &
\textbf{RS Origin} \\
\midrule
Strong (QCD) & $\theta_{\mathrm{eff}} = 0$ exactly &
$J$-symmetry $J(x) = J(x^{-1})$ for bare angle;
real $\varphi$-ladder masses for mass phase \\
Weak (CKM) & $\delta_{\mathrm{KM}} \neq 0$ &
Asymmetric 3-gen torsion $\{0, 11, 17\}$ with
chiral weak coupling \\
Neutrino (PMNS) & $\delta_{\mathrm{CP}} \approx -\pi/2$ &
8-tick quarter-period phase rotation \\
CPT & Exact &
Double-entry ledger + voxel locality + Lorentz
limit \\
\bottomrule
\end{tabular}
\end{center}

%% ============================================================
\section{Predictions and Falsifiers}
\label{sec:predictions}
%% ============================================================

\begin{enumerate}
  \item \textbf{Neutron EDM}: $d_n \lesssim
  10^{-32}\;e\cdot\mathrm{cm}$ (SM-only prediction).
  Detection of $d_n > 10^{-29}\;e\cdot\mathrm{cm}$ would
  indicate new CP violation beyond the SM and would require
  RS revision.

  \item \textbf{No QCD axion}: No axion coupling to
  $G\tilde{G}$ at any mass.  Confirmed detection of a QCD
  axion would falsify the RS resolution of the strong CP
  problem.

  \item \textbf{PMNS CP phase}:
  $\delta_{\mathrm{CP}} \approx 270^\circ$.
  DUNE and Hyper-Kamiokande will test this to $\pm 10^\circ$.

  \item \textbf{CPT invariance}: No CPT violation in any
  sector.  Observation of a particle--antiparticle mass
  difference (e.g., $m_K \neq m_{\bar{K}}$) would falsify the
  double-entry ledger.

  \item \textbf{Cabibbo angle}:
  $\lambda \approx \varphi^{-3} \approx 0.236$, within $5\%$ of the
  observed value $0.2253 \pm 0.0007$.  (Note: the generation-torsion
  step $\Delta\tau = 11$ gives $\varphi^{-11/2} \approx 0.071$, which
  does \emph{not} match $\lambda$; the correct RS identification is
  $\lambda \approx \varphi^{-3}$, as detailed in the Remark in
  Section~\ref{sec:weak-cp}.)
\end{enumerate}

%% ============================================================
\section{Conclusion}
\label{sec:conclusion}
%% ============================================================

The complete CP structure of physics is derived from RS:
\begin{enumerate}
  \item $\theta_{\mathrm{eff}} = 0$ from $J(x) = J(x^{-1})$
  and the reality of $\varphi$-ladder masses---no PQ axion
  needed.
  \item Weak CP violation from the asymmetric generation
  torsion $\{0, 11, 17\}$ acting on chiral weak doublets.
  \item PMNS CP phase $\delta_{\mathrm{CP}} \approx -\pi/2$
  from the 8-tick quarter-period.
  \item CPT exact from the double-entry ledger structure.
\end{enumerate}

The strong CP ``problem'' is dissolved: $\theta_{\mathrm{eff}} = 0$
is a structural identity in a framework where all masses are real
and the gauge action is built from the symmetric $J$-cost functional.

%% ============================================================
\begin{thebibliography}{99}

\bibitem{Aczel1966}
J.~Acz\'el, \emph{Lectures on Functional Equations and Their Applications},
Academic Press, New York (1966).

\bibitem{RS_Axioms}
J.~Washburn, ``The Algebra of Reality: A Recognition Science Derivation
of Physical Law,'' \emph{Axioms} \textbf{15}(2), 90 (2026).
\texttt{doi:10.3390/axioms15020090}

\bibitem{PQ}
R.~D.~Peccei and H.~R.~Quinn, ``CP Conservation in the
Presence of Instantons,'' Phys.\ Rev.\ Lett.\ \textbf{38},
1440 (1977).

\bibitem{PDG}
Particle Data Group, ``Review of Particle Physics,'' PTEP
\textbf{2022}, 083C01 (2022).

\bibitem{nEDM}
C.~Abel \textit{et al.}\ (nEDM Collaboration),
``Measurement of the Permanent Electric Dipole Moment of the
Neutron,'' Phys.\ Rev.\ Lett.\ \textbf{124}, 081803 (2020).

\bibitem{Jarlskog1985}
C.~Jarlskog, ``Commutator of the Quark Mass Matrices in the
Standard Electroweak Model and a Measure of Maximal CP
Nonconservation,'' Phys.\ Rev.\ Lett.\ \textbf{55}, 1039
(1985).

\bibitem{Pospelov2005}
M.~Pospelov and A.~Ritz, ``Electric Dipole Moments as
Probes of New Physics,'' Ann.\ Phys.\ \textbf{318}, 119
(2005).

\bibitem{Streater1964}
R.~F.~Streater and A.~S.~Wightman,
\textit{PCT, Spin and Statistics, and All That},
W.~A.~Benjamin (1964).

\bibitem{NuFIT}
I.~Esteban \textit{et al.}, ``The Fate of Hints: Updated
Global Analysis of Three-Flavor Neutrino Oscillations,''
JHEP \textbf{09}, 178 (2020).

\bibitem{DUNE}
DUNE Collaboration, ``Deep Underground Neutrino Experiment
(DUNE) Far Detector Technical Design Report,''
arXiv:2002.03005 (2020).

\bibitem{Sakharov1967}
A.~D.~Sakharov, ``Violation of CP Invariance, C Asymmetry,
and Baryon Asymmetry of the Universe,'' Pisma Zh.\ Eksp.\
Teor.\ Fiz.\ \textbf{5}, 32 (1967).

\bibitem{Planck2018}
Planck Collaboration, ``Planck 2018 Results. VI. Cosmological
Parameters,'' Astron.\ Astrophys.\ \textbf{641}, A6 (2020).

\end{thebibliography}

\end{document}
